\scanpage{13}\mspage{12}
\deleted{ce qui}

Composantes connexes de $M=$ ens des multistrates. Les éléments minimaux de $M$ sont ceux qui sont réduits à un seul ensemble $A$, \deleted{(qui peut être} $A\ne\emptyset$ (NB $\unsure{\mathscr{L}|\Phi}=\{\emptyset\}$ n'est pas une famille admissible !\deleted{)} Si $\mathscr{L}$ est admissible, $\emptyset\notin\deleted{\unsure{F}}\mathscr{L}$, \deleted{)} puisque $A\in\mathscr{L}\Rightarrow A^{\circ}\ne\emptyset$ a fortiori $A\ne\emptyset$)
\[
\mathscr{L}\deleted{\unsure{F_{A}|\mathscr{L}_{A}}}=\{A\}
\]

\deleted{(NB)} \deleted{(NB $F\in\unsure{F|M}\Rightarrow F\ne\emptyset$ !)}

Si $X\deleted{\unsure{\mathscr{L}|F}}\in M$, \deleted{et $A\in\deleted{\mathscr{L}}F$, alors \unsure{pour}} \inserted{$\deleted{\unsure{F}}X\subset\unsure{\mathcal P|\mathfrak P}(L)$}, et si $\deleted{F}X$ admet un \deleted{plus petit} élément minimal $A$ (p.ex. $\deleted{F}X$ fini, \inserted{on suppose $X\ne\emptyset$}) alors,
\[
\deleted{F}X_A=\{A\}\leq F
\]
\deleted{d'autre part,} donc $\deleted{F_A}$ \deleted{et} $\deleted{F}X$ et $\{A\}$ dans la $=$ comp. connexe. D'autre part, \deleted{\unsure{$\{A\}$ est}} si $A\ne L$, et si $x\in L\setminus A$, $\{A\}$ et $\{\{x\}\}$ sont dans la $=$ composante connexe, car ils sont majorés par le multistrate
\[
\bigl\{\{A\},\{\{x\}\},A\cup\{\{x\}\}\bigr\}
\]
\deleted{donc} Et pour la $=$ raison, \deleted{si} \deleted{Card} \unsure{on} voit que \inserted{$\{\{x\}\}$ pour tous les} \deleted{\unsure{tous les}} pts de $L$ sont dans la $=$ comp connexe de $L$. En résumé les $X\in M$ qui ont un plus petit élément \inserted{$A$} sont dans une $=$ composante connexe, du moment que $A\ne L$ i.e. du moment que $\unsure{X\prime|X}$ ne se réduise à $\{A\}=\{L\}$. Mais celui-ci aussi, si \deleted{Card} $L\deleted{\geq}\ne\deleted{\unsure{2|0}}\emptyset$, est dans la
